\scanpage{21}
tels termes. Donc $c_{i,j}\equiv0\ (p)$ dans ce cas, ce qui prouve (remplaçant $i,j$ par $u,v$) que $c_{u,v}\equiv0\ (p)$ donc encore $c_{ij}\equiv c_{u,v}\,c_{r,s}\ (p)$.

Le lemme 1 résulte du corollaire par une récurrence immédiate sur la longueur $e$ des développements $p$-adiques écrits de $u$ et de $v$.

\uline{Cor. 2} Soit $N$ un entier \inserted{$\geq0$} tel que $N+1\equiv0\ (p)$, posons $N+1=(N'+1)p$ donc $N=(N'+1)p-1=N'p+(p-1)$ \deleted{\unsure{$\in$|$($}} de sorte que à décalage de 1 près\unsure{\}} , \deleted{les} $N+1$ et $N'+1$ ont $=$ coeff. du développem\textsuperscript{t} $p$-adique, ainsi que $N$ et $N'$ (quand on fait abstraction du premier coeff. $p-1$ de $N$). Soit alors $\varepsilon\geq1$, on a
\[
\begin{aligned}
c_{N+1,p^\varepsilon-1}\deleted{\unsure{\not\equiv0\ (p)\iff}}&\equiv c_{N'+1,p^{\varepsilon-1}-1}\deleted{\unsure{\not\equiv0}}\ (p)\deleted{\unsure{\iff0}}\left\{\begin{aligned}&\equiv1\ (p)\ \text{si }N'+1\equiv0\ (p^{\varepsilon-1})\\&\equiv0\ (p)\ \text{si }N'+1\not\equiv0\ (p^{\varepsilon-1})\end{aligned}\right.\\
c_{N,p^\varepsilon}\deleted{\unsure{\not\equiv0\ (p)\iff}}&\equiv c_{N',p^{\varepsilon-1}}\deleted{\unsure{\not\equiv0}}\ (p)\deleted{\unsure{\iff}}\left\{\begin{aligned}&\equiv v_{\varepsilon-1}+1\ \inserted{(p)}\ \text{où}\\&\deleted{\unsure{v_\alpha}}\quad N'=\sum v_\alpha p^\alpha\end{aligned}\right)
\end{aligned}
\]

\uline{Dém.} Le développement $p$-adique de $p^\varepsilon-1$ est
\[
p^\varepsilon-1=(p-1)+(p-1)p+\cdots+(p-1)p^{\varepsilon-1}
\]
de m $p^{\varepsilon-1}-1=(p-1)+\cdots+(p-1)p^{\varepsilon-2}$ si $\varepsilon\geq2$ ($=0$ si $\varepsilon=1$).

Posant
\[
\left.\begin{aligned}
N'+1&=u_0+u_1p+\cdots+u_\nu p^\nu\\
N'&=v_0+v_1p+\cdots+v_\nu p^\nu
\end{aligned}\right\}\quad(\nu\geq\varepsilon-1)
\]
d'où
\[
\begin{aligned}
N+1&=0+u_0p+u_1p^{\unsure{2}}+\cdots+u_\nu p^{\nu+1}\\
N&=p-1+v_0p+v_1p^{\unsure{2}}+\cdots+v_\nu p^{\nu+1}.
\end{aligned}
\]
On trouve que
\[
\left\{\begin{aligned}
c_{N+1,p^\varepsilon-1}&\equiv c_{u_0,p-1}\,c_{u_1,p-1}\ \cdots\ c_{u_{\varepsilon-2},p-1}\quad(p)\qquad(\equiv1\ \inserted{(p)}\ \text{si }\varepsilon=1)\\
c_{N'+1,p^{\varepsilon-1}-1}&\equiv c_{u_0,p-1}\,c_{u_1,p-1}\ \cdots\ c_{u_{\varepsilon-2},p-1}\quad(p)\qquad(\equiv1\ (p)\ \text{si }\varepsilon=1)
\end{aligned}\right.
\]
\[
\left\{\begin{aligned}
c_{N,p^\varepsilon}&\equiv c_{v_{\varepsilon-1},1}\equiv v_{\varepsilon-1}+1\\
c_{N',p^{\varepsilon-1}}&\equiv c_{v_{\varepsilon-1},1}\equiv v_{\varepsilon-1}+1
\end{aligned}\right.
\]

\deleted{\unsure{Prop}}\uline{Lemme 2} Soit $\mathbb{Z}_p$ le localisé de $\mathbb{Z}$ en $p$, $E$ un $\mathbb{Z}_p$-module, $N\in\mathbb{N}$, $e\geq0$ tel que $p^e\geq N+1$ \deleted{\unsure{(donc $e\geq0$)}}, $I$ l'ens des entiers de la forme $N$ et $N+p^\varepsilon$, $0\leq\varepsilon\leq e$, $(x_i)_{i\in I}$ une famille d'éléments de $E$, tels qu'on ait
\[
\binom{j}{i}(x_j-x_i)=0\qquad i=N,N+1,\quad i<j\in I\quad(\text{donc }j=N+p^\varepsilon,\ 0\leq\varepsilon\leq e)
\]
\margin{\uline{Rmq} \unsure{même} si $E$ un $\mathbb{F}_p$-\unsure{es vectoriel|module}, on \unsure{ne} peut conclure que $x_i=x_N$ $\forall i\in I$ (ex : \inserted{$p=2$,} $N=2,e=2$, donc on a $x_2,x_3,x_4,x_6$ on a $x_2=x_3=x_6$ mais \uline{pas} néc. $=x_4$}% margin note written sideways
% archive sheet number: 20

